\cotefile{57}
\scanpage{14}
Carré, cube, rigidité\ldots

$\ast\ \unsure{\ast|\times}$

% archive sheet number: 13
\scanpage{15}
\[
\boxed{n^{*}(L)\simeq L^{\otimes n}(\delta L)^{\frac{n(n-1)}{2}}}
\]
\[
n+2\,\frac{n(n-1)}{2}=n^{2}
\]
\[
\delta L=L^{\otimes 2}\unsure{\delta'|\delta''}(L)\qquad
\delta'L_x=L_{2x}L_x^{-4}
\]
\[
\boxed{n^{*}(L)\simeq L^{\deleted{\unsure{n}}\unsure{n^{2}|2}\deleted{\unsure{-n}}}
(\delta'L)^{\frac{n(n-1)}{2}}}
\]

\diagram{\begin{tikzcd}[ampersand replacement=\&,column sep=small,row sep=small]
A\arrow[r]\& A\times A \& {} \arrow[d,phantom,"\unsure{\downarrow}"]\\
A\arrow[r]\&\widehat A\&{}
\end{tikzcd}}
\[
\boxed{\underline{\underline{\mathrm{Pic}}}^{\tau}_{(\prod X_i)/S}
\simeq \prod_i\underline{\underline{\mathrm{Pic}}}^{\tau}_{X_i/S}}
\]

\margin{\uuline{sans torsion}}
% slanted
% The product over unordered pairs, including its subscript, is circled.
\diagram{\begin{tikzcd}[ampersand replacement=\&,column sep=small,
execute at end picture={\draw (pairterm.center) ellipse[x radius=1.9cm,y radius=1.0cm];}]
0\arrow[r]\& \displaystyle\prod_i\underline{\underline{\mathrm{Pic}}}_{X_i/S}\arrow[r]\&
\underline{\underline{\mathrm{Pic}}}_{(\overset{\unsure{=|\rightleftarrows}}{\unsure{\prod_{\!S}|\prod}}X_i)/S}\arrow[r]\&
|[alias=pairterm]| \displaystyle\prod_{\substack{\{i,j\}\\\text{paires à}\\\text{deux}\\\text{éléments de}\\\text{\unsure{l'ens.} d'indices }I}}
\mathcal C_S(X_i,X_j)\arrow[r]\&0
\end{tikzcd}}

[\uline{\unsure{Splitté|Splitte}} si les $X_i/S$ ont tous \uuline{une section}]

Suite exacte de faisceaux [pour la descente fid. plate \unsure{qu. cpcte}]

% Circled, in the left margin.
\margin{\unsure{expliciter|explications} ce cas}
% slanted
les $X_i\longrightarrow S$ propres \inserted{et plats}, $\unsure{f_{*}(|f(}\mathcal O_{X_i})\simeq\mathcal O_S$ \unsure{universellt|universellement}.
% archive sheet number: 14
\scanpage{16}
\diagram{\begin{tikzcd}[ampersand replacement=\&,row sep=small]
X\arrow[d]\\ Y\ \deleted{\unsure{G|6}}
\end{tikzcd}}
\[
\mathrm{Pic}(X/Y)
\]

$\mathrm{Pic}_{\mathrm{sp}\deleted{\unsure{é}}}(X/Y)$ : classes à \unsure{isom|isom.} près des faisceaux inversibles $\underline L$, section marquée, telles que \deleted{\unsure{les} \unsure{div} soit} \deleted{(i)} l'on ait
\[
H^{i}_{f}\deleted{\unsure{(X,}}(\underline L\deleted{\unsure{(n))}})=0
\quad\text{pour }i>0,\ \deleted{\unsure{n\gg0}}
\]
donc $\underline H_f^0(\underline L)$ est loc libre de rang $\chi_{X/Y}(\underline L)$.
\deleted{$(ii)\quad\unsure{f_{*}}\ \unsure{\underline H^0_f}\ \chi_{X/Y}(\underline L)$}

% Wavy cancellation through the lower right-hand block.
\deleted{\unsure{on prend} $>1$, i.e l'\unsure{hom|hyp.} : \unsure{déduit|de div} de la section marquée $H^0_f(\underline L)\longrightarrow$}

\[
\underline H^0_f(\underline{\mathrm O}_X)=\underline{\mathrm O}_Y
\]
\deleted{\unsure{et}} ``universellement''

\unsure{une|ou} section marquée :

% This diagram is upside down relative to the definition.
\diagram{\begin{tikzcd}[ampersand replacement=\&,column sep=small,row sep=small]
\deleted{\unsure{P}}\&{}\&{}\\
\mathcal A\arrow[dd]\&\mathcal D\arrow[l]\&
\mathcal D\times_{\mathcal A}\mathcal D\arrow[l,shift left=.5ex]\arrow[l,shift right=.5ex]\\
{}\&{}\&\mathcal D\times_Y\mathcal D\arrow[u,phantom,"\cap"]\\
Y\& P\arrow[uu]\&
P\times_Y\deleted{\unsure{P}}\ \unsure{D_1D_2}\arrow[l,shift left=.5ex]\arrow[l,shift right=.5ex]\arrow[uu,bend right=16]
\end{tikzcd}}

Une telle donnée : \unsure{revient|équivaut} à la donnée
% At the foot only the tops of further strokes are visible.
$\unsure{F}\qquad\unsure{X}$
% archive sheet number: 15
\scanpage{17}
\diagram{\begin{tikzcd}[ampersand replacement=\&,column sep=small]
{}\&X_i\times X_j\arrow[r]\&\displaystyle\prod X_i\\
\deleted{X}X_{[ij\unsure{h|k}l]}\arrow[ur]\arrow[r]\&
X_h\times X_l\arrow[ur]\&{}
\end{tikzcd}}
\[
X_{ijhl}=\begin{cases}
X_{\varnothing}&\text{si }\{i,j\}\cap\{h,l\}=\varnothing\\
X_{\alpha}&\text{si }\{i,j\}\cap\{h,l\}=\{\alpha\}
\end{cases}
\]
\margin{$x_i\,x_j\quad x_h\,x_l$}

\figure{\begin{tikzpicture}[x=1cm,y=1cm]\draw[line width=1.4pt] (0,.09)--(0,0)--(1.65,0)--(1.65,.11) (1.4,0)--(1.4,.10);\end{tikzpicture}}

\[
L_{m_1x_1+\cdots+m_n x_n}
\]
\[
L_{x,y}^{\deleted{\unsure{2}}}\simeq
L_{xe}^{\deleted{\unsure{2}}}L_{ye}^{\deleted{\unsure{2}}}\mathcal M_{xy}
\]
\[
\mathcal M_{xy}\ \mathcal M_{yx}
\]
\[
\deleted{L_{xy}L_{yz}}
\]

\margin{$\begin{aligned}
&\deleted{n^2-3n+2-n^2+n}\\
&\deleted{=-2(n\unsure{+|-}1)}\\
1&=-\frac{n(n-1)}{2}+n+\frac{(n-1)(n-2)}{2}\\
\deleted{2}&=-\deleted{n^2}+\deleted{n}+\deleted{2n}+\deleted{n^2}-\deleted{3n}+\deleted{2}
\end{aligned}$}

\[
L_{x_1+\cdots+x_n}=\left(\prod_{i<j}L_{x_i+x_j}\right)
\left(\prod L_{x_i}^{-n+2}\right)L_{\varnothing}^{\alpha_n}
\]
\[
\deleted{L_{x+y+z}}
\]
% The last factor on this line is circled.
\[
L_{x+y}=L_xL_y\,\delta L_{xy}^{\deleted{\unsure{2}}}
\]
\margin{symétrique}
\margin{$(\delta L)_{0,0}=L_{\varnothing}^{-1}$}
\[
\prod_{i\nej} L_{x_i+x_j}=
\left(\prod L_{x_i}^{n-1}\right)\prod_{i\ne j}
(\delta L\deleted{\unsure{^{2}}})_{x_i,x_j}
\]
\[
\boxed{L_{\unsure{x_1\cdots x_n|x_1+\cdots+x_n}}=
\deleted{\unsure{\prod}}\prod_{i\unsure{<|\le}j}\delta L_{x_i,x_j}\,
\deleted{\unsure{L_{\varnothing}^{\alpha_n}}}\prod_i L_{x_i}\,L_{\varnothing}^{\alpha_n}}
\]
\[
\text{si }x_i=0\quad\text{\unsure{pour tt } }i\qquad
\underbrace{L_{\varnothing}^{-\frac{n(n-1)}2}
L_{\varnothing}^{\frac{(n-1)(n\unsure{-1|-2})}2}}_{L_{\varnothing}^{-n+1}}
\qquad \deleted{\unsure{(n-1)(n-1)}}
\]
% archive sheet number: 16
\scanpage{18}
\[
\begin{aligned}
1+1-1&=1\\
-(n-2)-(n-\unsure{2|7})+(n-3)&=-(n-1)\\
\alpha_n+\alpha_n-\alpha_{n-1}+1&=\deleted{\unsure{\alpha_n+\alpha_{n-1}}}\,2\alpha_n-\alpha_{n-1}+1=\alpha_{n+1}
\end{aligned}
\]
\[
\begin{aligned}
\alpha_3&=1\\ \alpha_4&=3\\ \alpha_5&=6\\ \alpha_6&=10\\ \alpha_7&=15
\end{aligned}
\]
% One cross through this four-line trial.
\[
\begin{aligned}
\deleted{\alpha_n=an+b}&\\
\deleted{2\alpha_n-\alpha_{n-1}+1}&\deleted{=2an+2b-an+a-b+1}\\
&\deleted{=an+b\unsure{+a|+b}+1)}\\
\deleted{\alpha_{n+1}}&\deleted{=an+a+b}
\end{aligned}
\]
\[
\begin{aligned}
\alpha_n&=an^2+bn+c\\
2\alpha_n-\alpha_{n-1}+1&=2an^2+2bn+2c-a(n^2-2n+1)-b(n-1)-c\inserted{+1}\\
&=an^2+(b+2a)n+(\deleted{\unsure{2c|a}}c-a+b+1)\unsure{-}\\
\alpha_{n+1}&=a(n^2+2n+1)+b(n+1)+c\\
&=an^2+(2a+b)n+(a+b+c)
\end{aligned}
\]
\[
\begin{aligned}
a+b+c&=-a+b+c+1\\
a&=-a+1\qquad a=\tfrac12
\end{aligned}
\]
\[
\alpha_n=\tfrac12 n^2+bn+c
\]
\[
\begin{aligned}
\alpha_2&=0\qquad \deleted{\unsure{c=0}}\qquad &\tfrac12\,4+2b+c&=0\\
\alpha_3&=1&&\tfrac12\,9+3b+c=1
\end{aligned}
\]
\[
\begin{array}{rl}
2b+c&=-2\\6b+2c&=-7\\ \hline
2b&=-\unsure{3|5}\\c&=+1
\end{array}
\]
\[
\begin{aligned}
&\tfrac12 n^2-\tfrac32 n+1\\
&\tfrac12[n^2-3n+2]
\end{aligned}
\]
% The fraction within the box is also circled.
\[
\boxed{\alpha_n=\frac{(n-2)(n-1)}{2}}
\]
\[
\deleted{L_{x_1\ldots x_n}\simeq
\left(\prod_{i\ne j}L_{x_ix_j}\right)
\left(\prod_{\unsure{\!1|i}} L_{x_i}^{-1}\right)L_{\varnothing}^{\frac{(n-1)\unsure{(n-2)}}{2}}}
\]
\[
\boxed{L_{x_1\ldots x_n}\simeq \left(\prod_{i\ne j}L_{x_ix_j}\right)
\left(\prod_i L_{x_i}\right)^{-n+2}L_{\varnothing}^{\frac{(n-1)(n-2)}2}}
\]
% archive sheet number: 17
\scanpage{19}
\diagram{\(\begin{array}{ccc}x&y&z\\a&b&c\end{array}\)}
\[
\boxed{L_{x,y,z}\simeq L_{xy}L_{y,z}L_{x,z}L_x^{-1}L_y^{-1}L_z^{-1}L_{\varnothing}}
\]
\margin{``\unsure{Th|Th.} du cube''\par\ldots}

\[
\begin{aligned}
L_{\underbrace{xy}zt}
 &=L_{(xy)z}L_{z,t}L_{(xy),t}L_{(xy)}^{-1}L_z^{-1}L_t^{-1}L_{\varnothing}\\
 &=\deleted{\unsure{L}}\underset{|}{L_{xy}}\underset{|}{L_{y,z}}\underset{|}{L_{xz}}
 \underset{|}{L_x^{-1}}\underset{|}{L_y^{-1}}\underset{|}{L_z^{-1}}\inserted{L_{\varnothing}}
 \underset{|}{L_{zt}}\deleted{L_{xy}}\underset{|}{L_{y,t}}\underset{|}{L_{x,t}}
 \underset{|}{L_x^{-1}}\underset{|}{L_y^{-1}}\underset{|}{L_t^{-1}}\inserted{L_{\varnothing}}
 \deleted{L_{xy}^{-1}}\underset{|}{L_z^{-1}}\underset{|}{L_t^{-1}}\underset{|}{L_{\varnothing}}
\end{aligned}
\]
\[
\boxed{L_{xyzt}=L_{xy}^{\deleted{\unsure{2}}}L_{xz}L_{xt}L_{yz}L_{yt}L_{z,t}
\deleted{\unsure{L}}L_x^{-2}L_y^{-2}L_z^{-2}L_t^{-2}L_{\varnothing}^{3}}
\]
\[
\begin{aligned}
L_{\underbrace{xyz}\underbrace{tu}}&=L_{xyzt}L_{tu}L_{xyzu}L_{xyz}^{-1}L_t^{-1}L_u^{-1}L_{\varnothing}\\
&=\deleted{L_{xy}}\underset{/}{L_{xz}}\underset{/}{L_{xt}}\underset{/}{L_{yz}}\underset{/}{L_{yt}}\underset{/}{L_{zt}}\underset{/}{L_x^{-2}}\underset{/}{L_y^{-2}}\underset{/}{L_z^{-2}}\underset{/}{L_t^{-2}}\underset{/}{L_{\varnothing}^3}\underset{/}{L_{tu}}\deleted{\unsure{L}}\\
&\quad \underset{/}{L_{xy}}\deleted{L_{xz}}\underset{/}{L_{xu}}\deleted{L_{yz}}\underset{/}{L_{yu}}\underset{/}{L_{zu}}\underset{/}{L_x^{-2}}\underset{/}{L_y^{-2}}\underset{/}{L_z^{-2}}\underset{/}{L_u^{-2}}\underset{/}{L_{\varnothing}^3}\\
&\qquad \deleted{L_{xy}^{-1}}\deleted{L_{yz}^{-1}}\deleted{L_{xz}^{-1}}\underset{/}{L_x}\underset{/}{L_y}\underset{/}{L_z}
\underbrace{\deleted{L_{\varnothing}^{-1}}\underset{/}{L_t^{-1}}\underset{/}{L_u^{-1}}\deleted{L_{\varnothing}}}
\end{aligned}
\]
\[
\boxed{\begin{aligned}
L_{xyztu}&=L_{xy}L_{xz}L_{xt}L_{xu}L_{yz}L_{yt}L_{yu}L_{zt}L_{zu}L_{tu}\\
&\quad L_x^{-3}L_y^{-3}L_z^{-3}L_t^{-3}L_u^{-3}\deleted{\unsure{1}}\\
&\qquad L_{\varnothing}^{\unsure{6|0}}
\end{aligned}}\qquad L_{\varnothing}^{6}
\]
\[
L_{x_1\ldots x_n}=\prod_{i<j}L_{x_i x_j}\ \prod L_{x_i}^{-(n-2)}\,L_{\varnothing}^{\alpha_n}
\]
\[
\begin{aligned}
% The two long return strokes lead from each right-hand line end to the following line.
L_{x_1\ldots x_{n+1}}&=L_{(x_1\ldots x_{n\unsure{-2|-1}})x_n x_{n+1}}=
 L_{x_1\ldots x_n}L_{x_1\ldots x_{n-1}x_{n+1}}L_{x_nx_{n+1}}
 L_{x_1\ldots x_{n-1}}^{-1}L_{x_n}^{-1}L_{x_{n+1}}^{-1}L_{\varnothing}\\
 &=\prod_{\substack{\deleted{\unsure{i<j}}\\1\le i<j\le n}}L_{x_i x_j}
 \prod_{1\le i\le n}L_{x_i}^{-(n-2)}L_{\varnothing}^{\alpha_n}
 \prod_{\substack{\deleted{\unsure{i<j}}\\1\le i<j\le n+1\\i,j\ne n}}L_{x_i x_j}
 \prod_{\substack{1\le i,j\le\unsure{n-1|n+1}\\i,j\ne n}}L_{x_i}^{-(n-2)}L_{\varnothing}^{\alpha_n}\\
 &\quad L_{x_n\deleted{\unsure{2}},x_{n+1}}
 \prod_{\substack{\deleted{\unsure{1\le i\le n-1}}\\\deleted{\unsure{i<j}}\\1\le i<j\le n-1}}L_{x_ix_j}^{-1}
 \prod_{1\le i\le n-1}L_{x_i}^{+(n-3)}L_{\varnothing}^{-\alpha_{n-1}}\\
 &\qquad L_{x_n}^{-1}L_{x_{n+1}}^{-1}L_{\varnothing}
\end{aligned}
\]
\[
\begin{aligned}
&1\le i,j\le n-1\\
&i\ne j
\end{aligned}
\]
\[
\left(\prod_{\substack{\deleted{\unsure{i\ne j}}\\\deleted{\unsure{i<j}}}}L_{x_ix_j}\right)
\left(\prod L_{x_i}^{-(n-1)}\right)L_{\varnothing}
\]
% archive sheet number: 18
\scanpage{24}
Rigidité

Cf \unsure{aussi} $\overline{\underline{\mathrm{V}}}$
% archive sheet number: 23
\scanpage{25}
RIGIDITÉ ET GROUPES ALG. \mspage{1}

\margin{``Theorem of the square''}
% slanted
\diagram{\begin{tikzcd}[ampersand replacement=\&,column sep=small,row sep=small]
{}\& X_1\times_S X_2\arrow[dl,"p_1"']\arrow[dr,"p_2"]\&{}\\
X_1\arrow[dr,"\varphi_1"',dash]\&{}\&X_2\arrow[dl,"\varphi_2",dash]\\
{}\&S\&{}
\end{tikzcd}}

\uline{Théorème} Soient $S$ un préschéma \inserted{\unsure{loc. noeth.}}, $\varphi_i:X_i\longrightarrow S$ $(i=1,2)$ deux préschémas de type fini sur $S$, \inserted{$\varepsilon_1$ une section de $X_1$ sur $S$} \deleted{$G$ un schéma en groupes loc. de type fini sur $S$, on suppose} soient $p_i$ \inserted{$(i=1,2)$} les deux projections de $X_1\times_S X_2$ sur les facteurs. On suppose a) $p_1$ est propre, et $p_{1*}(\underline{\mathrm O}_{X_1\times X_2})\xleftarrow{\sim}\underline{\mathrm O}_{X_1}$

\deleted{\unsure{b)} Pour tout voisinage ouvert $U$ de $\varepsilon_1(S)$, \unsure{$X_1$ est} \illegible{1} $U\times_S X_2$ est schématiquement dense dans $X_1\times_S X_2$}

\inserted{les fibres géom de $X_1/S$ sont connexes}
\deleted{\unsure{a)}} b) $p_1$ est universellement ouvert, \inserted{\uuline{ou} $X_1$ plat sur $S$, à fibres géom primaires}

[ \uline{Ex} $\varphi_2$ est propre, plat et $s\in S\Rightarrow k(s)\simeq H^0(X_{1s},\underline{\mathrm O}_{X_{1s}})$ ; \deleted{``\unsure{d'ailleurs|primaires}'' \unsure{les sous-algèbres art. intègres, etc|les sous-algèbres art. intègres|i.e. sans cycles ass. immergés}} $\varphi_1$ \deleted{est plat,} à fibres géom.$^{\mathrm t}$ irréductibles ]. Alors pour tt préschéma en groupes $G$ sur $S$, séparé sur $S$, et tt $S$-hom
\[
f:X_1\times_S X_2\longrightarrow G
\]
on a
\[
f=(f_1p_{\deleted{\unsure{2}}1})(\unsure{f_2|f}p_2)\tag*{(*)}
\]
où
\[
f_1:X_1\longrightarrow G,\qquad f_2:X_2\longrightarrow X
\]
sont des $S$-homomorphismes, déterminés de façon unique par la condition supplémentaire
\[
\deleted{f_2=f\circ(\varepsilon_1\times_S\operatorname{id}_{X_2})}
\]
\deleted{[ équivaut à :} $(**)$ \quad $f_1\varepsilon_1={}$ section unité, \deleted{pourvu que $\varphi_2$ soit un épim. ]} \inserted{et alors $(***)\ f_2=f\circ(\varepsilon_1\times\operatorname{id}_{X_2})$}

\margin{Attention ! on ne suppose pas la loi de composition commutative}
% slanted
\margin{plus haut}
% slanted
\uline{Corollaire} Supposons de plus que $X_2$ soit muni d'une section $\varepsilon_2$ sur $S$. Alors $f$ se met de façon unique sous la forme
\[
f_1p_1+f_2p_2+\sigma\varphi
\]
\inserted{tels que $f_i\varepsilon_i={}$ section unité;} où $f_i:X_i\longrightarrow G$ sont des $S$-morphismes, $\sigma:S\longrightarrow G$ une section de $G$ sur $S$, [ $\varphi$ désignant \deleted{$\varphi_1\times_S\varphi_2$} la projection $X_1\times_S X_2\longrightarrow S$ ].
% archive sheet number: 24
\scanpage{26}
Démonstration du th.

Soit
\[
f_2=f\circ(\varepsilon_1\times_S\operatorname{id}_{X_2})
\]
et considérons
\[
f'=f-f_2p_2:X_1\times_S X_2\longrightarrow G
\]
\deleted{En vertu de a) et} \deleted{\unsure{Par}} On a donc $f'\circ(\varepsilon_1\times\operatorname{id}_{X_2})={}$ morphisme unité $X_2\to G$, \deleted{i.e.} \inserted{donc} pour tout $x_1\in\varepsilon_1(S)$, le morphisme induit par $f':k(x_1)\otimes_{k(s)}X_2\longrightarrow G$ \deleted{est} a une image finie. En vertu de l'hyp. a) et du lemme 1 ci-dessous, il en résulte qu'il existe un voisinage $U_1$ de $\varepsilon_1(S)$ telle que $f'|U_1\times_S X_2$ soit de la forme $f_1p_1$, où $f_1:U_1\longrightarrow G$ \unsure{un} $S$-morphisme convenable, donc pour tout $x_1\in U_1$, $f'$ applique $k(x_1)\otimes_S X_2$ en une partie finie de $G$. En vertu du lemme 2 ci-dessous, \inserted{et $p_1$ étant univ. ouvert \deleted{\unsure{puisque}} par b),} cela reste vrai pour tt $x\in\overline U_1$, i.e. $\forall\,x\in X_1\times_S X_2$ grâce à c). En vertu du lemme 1, on peut donc prendre $U_1=X_1$.

\margin{\diagram{\begin{tikzcd}[ampersand replacement=\&,column sep=tiny,row sep=tiny]
X_1\times_S X_2\arrow[d]\arrow[dr,bend left]\&{}\\
X_1\&X_1\times G\arrow[l]
\end{tikzcd}}}
Cela prouve l'existence de $f_1,f_2$ ayant les propriétés \deleted{\unsure{dites}} \inserted{$(*)$}. \deleted{[ $f_1\varepsilon_1={}$ section unité est évident ].} la propriété $(**)$ s'en déduit par une vérification triviale \unsure{en} \deleted{de la définition de $f_2$, \unsure{bien}} [ \inserted{remplacer $f_1$ par $f_1f_1^{\prime-1}$.} \deleted{qui faite \unsure{de} $f= f_1p_1+f_2p_2$ et \unsure{$f_2$}} \inserted{où $f'_1=f_1\varepsilon_1$} \deleted{d'où \unsure{en composant} au premier membre \unsure{relation} avec $\varepsilon_1\times\operatorname{id}_{X_2}$ : $(f_1\varepsilon_1)\varphi_2={}$ \unsure{section} unité} \inserted{morphisme unité}, car on voit tout de suite que \inserted{D'ailleurs,} \deleted{donc} On a, en supposant $f_2=f\circ(\varepsilon_1\times\operatorname{id}_{X_2})$
% archive sheet number: 25
\scanpage{27}
\mspage{2}
\uline{Lemme 1} Soient $f:X\longrightarrow Y$ un morphisme propre avec $Y$ loc. \unsure{noeth|noeth.}, tel que $f_*(\underline{\mathrm O}_X)\xleftarrow{\sim}\underline{\mathrm O}_Y$, \deleted{et soit} $\deleted{\unsure{g}}h:Z\longrightarrow\deleted{\unsure{X}}\unsure{Y}$ un \deleted{$Y$-préschéma de type fini,} \inserted{morphisme séparé loc. de type fini,} $g:X\longrightarrow Y$ un $Y$-morphisme, \deleted{supposons que} $y\in Y$, supposons que $g(X_y)\subset Z_y$ soit \uline{fini}. Alors il existe un voisinage \inserted{ouvert} $U$ de $y$ et une section $\sigma$ de $Z$ sur $U$, tels que $g|f^{-1}(U)=\sigma\circ(f|f^{-1}(U))$. Pour $U$ donné, $\sigma$ est unique.

\margin{\unsure{rédiger ceci|rédiger en} en \unsure{énonçant|invoquant|important} sous forme du corollaire plus bas}
% slanted

\uline{Dém.} Considérons $g(X)=X'\subset Z$, il est propre sur $Y$, et sa fibre $X'_y$ en $y$ est finie, donc $\exists\ U\ni y$ tq $X'|U$ soit \uline{fini} sur $U$. On peut supposer $U=Y$. Soit $\mathcal A=g_*(\underline{\mathrm O}_X)$, c'est un faisceau cohérent d'Algèbres sur $Z$, définissant un $Z'$ fini sur $Z$, et en fait fini sur $Y$, et on a une factorisation canonique $g=g''g'$. On a

\margin{\diagram{\begin{tikzcd}[ampersand replacement=\&,row sep=tiny,column sep=tiny]
{}\&Z'\arrow[d,"g''"]\\
X\arrow[ur,dashed,"g'"]\arrow[r,"g"']\arrow[d,"f"']\&Z\arrow[dl,"h"]\\
Y\&{}
\end{tikzcd}}}
\[
\deleted{\unsure{g''_{*}(\underline{\mathrm O}|g''(\underline{\mathrm O}}}\,
\underline{\mathrm O}_{Y}\longrightarrow h_*(\mathcal A)
\subset f_*(\underline{\mathrm O}_X)\simeq\underline{\mathrm O}_Y,
\qquad\inserted{\mathcal A\subset g_*(\underline{\mathrm O}_X)}
\]
d'où $(h\circ g'')_*(\underline{\mathrm O}_{Z'})\simeq\underline{\mathrm O}_Y$, i.e. $Z'\longrightarrow Y$ est un \uline{isom.} On prend alors $\sigma=g''(hg'')^{-1}$. L'unicité \deleted{\unsure{de}} de $\sigma$ est évidente. De cette unicité résulte

\uline{Corollaire} Sous les conditions du lemme, l'ens $U$ des $y\in Y$ satisfaisant la condition $g(X_y)$ fini est ouvert, et il existe une section unique
% archive sheet number: 26
\scanpage{28}
$\sigma$ de $Z$ sur $U$ telle que $g|f^{-1}(U)=\sigma(f|f^{-1}(U))$

\uline{Lemme 2} Soient $f:X\longrightarrow Y$ un morphisme de type \uline{fini} \uline{universellement} \uline{ouvert}, \deleted{alors l'ens des pts de} $h:Z\longrightarrow Y$ un morphisme \inserted{loc} de type fini, $h:X\longrightarrow Z$ un $Y$-morphisme, alors l'ens $U$ des $y\in Y$ tels que $g(X_y)$ soit \uline{fini} est \deleted{\uline{ouvert}} \uline{fermé}.

\inserted{On peut supposer $Y$ noethérien, \deleted{\unsure{affine}} $X$ et $Z$ aussi.} Il est \unsure{évidemt} constructible, donc on est ramené au cas où $Y$ est le spectre d'un anneau de valuation discrète, et où le pt générique $\unsure{\eta|y}$ est dans $U$. On peut supposer $X$ irréductible, \deleted{donc} \inserted{il} domine nécessairement $Y$, alors $\overline{g(X)}$ est irréductible et domine $Y$, \deleted{\unsure{sa}} la fibre générique est finie, donc aussi la fibre spéciale, OK.
\margin{\unsure{$y$}}

\deleted{\uline{Corollaire 1} Supposons $p_{1*}(\underline{\mathrm O}_{X_1}$}

\uline{Remarque.} \deleted{Le lemme 2 est} \inserted{Dans le corollaire 1 il suffit de} \deleted{\unsure{la condition en}} \inserted{simplement} supposer que \deleted{\unsure{l'on} \unsure{a dans le corollaire le}} \inserted{$p:X_1\times_S X_2\longrightarrow S$ est propre et $p_*(\underline{\mathrm O}_{X_1\times X_2})\xleftarrow{\sim}\underline{\mathrm O}_S$} \deleted{\unsure{Appliqu|quasi} \unsure{sur chaque} \unsure{fibre},
$\unsure{f'_2|f_2}=f-f\circ\varepsilon_1p_1-f\circ\varepsilon_2$,} En effet, \unsure{soit|on a} $f'=f-f_1p_1-f_2p_2$ $(f_i=f\circu_i\ \ldots)$, il reste à prouver que $f'(\text{\unsure{fibre|fibres}})$ est fini, pour ceci on \unsure{aura à se placer} sur un corps et appliquer lemme 2.
% archive sheet number: 27
\scanpage{29}
\mspage{3}
\uline{Corollaire} 2 Soit $\deleted{\unsure{X|A}}A$ un schéma abélien sur $S$ loc. noethérien, et $G$ un préschéma en groupes loc. de type fini sur $S$. Alors \deleted{\unsure{tout}} \deleted{\unsure{$S$-}}hom \inserted{de $S$-préschémas} de $A$ dans $G$ est de la forme (de façon unique)
\[
f=u+c
\]
\margin{Attention, si $G$ n'est pas commutatif}
% slanted
où $u$ est un hom en groupes, et où $c$ est un hom ``constant'' \inserted{i.e. provient d'une section de $G$ sur $S$}. On a nécessairement $c=f\circ\varepsilon_A$, où $\varepsilon_A$ est la section unité de $A$. \unsure{Donc} si $f\circ\varepsilon_A=\varepsilon_G$, $f$ est un hom de schémas en groupes.

\uline{Corollaire} 3 Soit $A$ un schéma abélien sur $S$. Alors $A$ est commutatif. Toute loi de composition sur $A$, telle que $\varepsilon_A$ \deleted{\unsure{soit}} soit unité bilatère, est la loi de composition donnée de $A$.

Il suffit de prouver cette dernière assertion. Or on a, si $f:A\times A\longrightarrow A$ est la loi de composition, en vertu du cor. 1 :
\[
f=[f\circ(\varepsilon\times\operatorname{id}_A)]p_1+
(f\circ(\operatorname{id}_A\times\varepsilon))p_2+
\deleted{\unsure{f\circ(\varepsilon_1}}[f\circ(\varepsilon_1\times\varepsilon_2)]p_2,
\]
i.e. $f=p_1+p_2$, c'est la loi de composition donnée. OK,
% archive sheet number: 28
\scanpage{30}
\uline{Reformuler le th.}

\uline{Th} $X$ propre et plat sur $S$, $H^0(X_s,\underline{\mathrm O}_{X_s})\xleftarrow{\sim}k(s)$ pour tt $s\in S$, ($S$ loc. \unsure{nœth|noeth.}), $\varepsilon_X$ une section de $X/S$, \deleted{\unsure{$\mathrm{Hom}(X$|$\mathrm{Hom}(X,$}} $G$ un foncteur sur $(\unsure{\mathrm{Sch}|\mathrm M})_{/S}^{\circ}$ à valeurs dans la catégorie des groupes. \deleted{On suppose donné un} \inserted{$f:X\longrightarrow G$ un $S$-hom pointé.} \deleted{On suppose donné un hom}

(i) Si $S_0$ est un sous-préschéma \deleted{\uline{\unsure{réduit}}} de $S$ ayant $=$ ensemble sous-jacent, et si \deleted{$G$ est} l'hom $f_0:X_0\longrightarrow G_0$ est trivial, alors $f$ l'est [ pourvu que \unsure{--- petite condition|--- petite cartésienne|--- petite catégorie} )

(ii) si de plus $G$ commute avec $\varprojlim$ d'anneaux artiniens, aux $\varinjlim$ d'anneaux noeth. \inserted{[cas de la localisation]}, et \deleted{\unsure{les}} les morphismes $\operatorname{Spec}(\widehat A)\longrightarrow\operatorname{Spec}(A)$ sont des morphismes de $G$-\unsure{0|O}-descente, \uline{alors} si $f_{\unsure{\times s|s}}=f\times_S k(s)$ est trivial, $f$ est trivial dans un voisinage

(iii) Si $G$ est séparé sur $S$, alors sous les conditions de (ii), \deleted{\unsure{on}} $f$ est trivial sur la composante connexe de $\unsure{e|s}$

En d'autres termes, $\underline{\mathrm{Hom}}_S^{\bullet}(X,G)$ est \uline{non ramifié} sur $\unsure{S|?}$
% archive sheet number: 29
\scanpage{31}
% archive sheet number: 30
\uline{Th}

$f:X\longrightarrow S$ propre et plat, $\underline O_S\simeq f_*(\underline O_X)$ universellement, $S$ loc. noeth.

$G$ foncteur $\mathrm{Sch}_{/S}^{o}\longrightarrow(\mathrm{Ens})$

% slanted
\margin{Utiliser le cas des \unsure{immersions|immersions ouvertes} -----}

a. $G$ formellement représentable ---

% slanted
\margin{Par \unsure{ailleurs|ex.|exemple}, il faut que $\unsure{\mathrm{Sch}'_{G/S}|\mathrm{Sch}_{G/S}}$ soit \unsure{rel$^{\mathrm{t}}$|un} représentable}

\unsure{b|d}) $G$ ``compatible avec \deleted{$\varprojlim$ d'anneaux} \unsure{Gogo|Cogo}''

b) $G$ \inserted{de nature locale} compatible avec descente fid. plate et qu. \unsure{cpcte|cpt}

\unsure{a|c}) $G$ compatible avec $\unsure{\varinjlim|\underline{\lim}}$ \deleted{d'\unsure{extensions} schémas} \inserted{d'anneaux}

% margin note written sideways
\margin{Conditions vérifiées si $G$ est représentable loc. de type fini sur $S$}

$u:X\to G$ i.e. $u\in G(X)$

Soit $U$ l'ens des $s\in S$ tels que $u_s\in G(X_s)$ soit l'image d'un élément de $G(k(s))$, i.e. $u_s:X_s\to G$ se factorise en
\[
X_s\longrightarrow k(s)\longrightarrow G
\]

Alors \quad 1) $U$ est ouvert [et \uline{fermé} si on suppose ----- \inserted{(critère valuatif)}]

2) Il existe un \deleted{morphisme} $v\in G(U)$ \unsure{unique} dont l'image dans $G(X|U)$ soit $u|(X|U)$, \deleted{i.e.}

\uline{Corollaire 1} Si $S$ est connexe, et s'il existe un $s\in S$ tel que $u_s:X_s\to G$ se factorise par $k(s)$, alors $u:X\to G$ se factorise par $X\xrightarrow{f}S\xrightarrow{v}G$

\uline{Corollaire 2} Supposons que $X$ et $G$ soient munis de sections sur $S$, et que $u$ soit compatible avec les sections marquées. Si $S$ est connexe, et si pour \uline{un} $s\in S$, $u_s$ est trivial, alors $u$ est trivial

\uline{Corollaire} 3 \quad $X$ sur $S$, \deleted{\unsure{muni}} $G$ sur $S$ comme \deleted{dessus} \inserted{au cor. 2} \deleted{\unsure{$X$|$Y$} munis de sections}, $Y$, un $S$ loc. de type fini sur $S$, \inserted{$Y$ connexe,} et $u:X\times_S Y\to G$ un hom. induisant \deleted{l'identité sur} l'hom trivial sur $\unsure{\varepsilon|e}_X\times_S Y$, \deleted{et} S'il existe $y\in Y$ tel que $u_y:X\times k(y)\to G$ soit trivial, alors $u$ est trivial.
\scanpage{32}
% archive sheet number: 31
\uline{Cor.} 4 \quad Sous les conditions du cor 2, soit de plus $Y$ loc. de type fini sur $S$, \inserted{à fibres géom. connexes, $Y$} muni d'une section $\unsure{\varepsilon|e}_Y$, Soit $u:X\times Y\to G$ induisant les morphismes triviaux $\varepsilon_X\times Y\to G$ et $X\times\varepsilon_Y\to G$, alors $u$ est trivial.

\uline{Cor.} 5 \inserted{(Weil)} Sous les conditions du cor. 4, supposons de plus $G$ un \deleted{foncteur en groupes, la section}

Alors tout morphisme $u:X\times_S Y\to G$ se met de façon unique sous la forme
\[
u=u_1\mathrm{pr}_1+u_2\mathrm{pr}_2+v\deleted{\unsure{p|p_S}}
\]
$\mathrm{pr}_i$ les deux projections $X\times_S Y\substack{\longrightarrow X\\\longrightarrow Y}$, $p:X\times_S Y\to S$ le morphisme structural, $u_1:X\to G$, $u_2:X\to G$, $v:S\to G$ des $S$-morphismes convenables.

[Demander Mumford]

\uline{Question} $S$ noeth (on est ramené au cas : $S$ spectre d'un anneau de val discrète \unsure{disons|sinon|géom.} complet), $X,Y$ propres et plats sur $S$, tels que $f_*(\mathcal O_X)\xleftarrow{\sim}\mathcal O_S$ et $g_*(\mathcal O_Y)\xleftarrow{\sim}\mathcal O_S$ univ$^{\mathrm t}$. $X$ et $Y$ munis de sections $e_X$ et $e_Y$. Soit $\underline L$ Module inversible sur $X$, \deleted{\unsure{avec}} trivialisations données sur $e_X\times_S Y$ et $X\times_S e_Y$, [compatibles sur $e_X\times_S e_Y$], et supposons qu'il existe un ouvert dense $U$ tel que $\underline L|\deleted{\unsure{X|X\times}}X\times_S Y|U$ soit trivial. Alors $\underline L$ est trivial

[En d'autres termes, $\underline{\underline{\mathrm{Corr}}}_S(X,Y)$ est tjrs \uline{séparé sur $S$} ?] \unsure{--- même|; ---|; c'est vrai} si \deleted{\unsure{$X$|$X$ et}} $\underline{\underline{\mathrm{Pic}}}_{X\times_S Y}$ ne l'est pas, par \deleted{\unsure{contre|exemple}} si $\underline{\underline{\mathrm{Pic}}}_{X/S}$ et $\underline{\underline{\mathrm{Pic}}}_{Y/S}$ ne le sont pas ]

Réponse affirmative si $\underline{\mathrm{Pic}}_{X/S}$ \uline{ou} $\underline{\mathrm{Pic}}_{Y/S}$ existe \inserted{\unsure{sur les anneaux de val discrète complets|sur les anneaux de val. discrète complets}}. [\unsure{écrire|voir}] \deleted{[p.ex. si $f$ ou $g$ à fibres séparables]}
% slanted
\margin{Il suffit \unsure{=|---} que $\underline{\mathrm{Pic}}_{X/S}$ ou $\underline{\mathrm{Pic}}_{Y/S}$ soit \unsure{pré-ét-représentable|pré-êt-représentable}, ce qui est le cas si $X$ \uline{ou} $Y$ est \deleted{\unsure{propre}} projectif sur $S$}
% slanted
\margin{\deleted{\unsure{ou|alors}} si $f$ \uline{ou} $g$ a un foncteur de Picard \uline{séparé} sur $S$.}
\scanpage{33}
% archive sheet number: 32
\inserted{(Déformation de schémas abéliens)}

\uline{Théorème 1} \quad $\deleted{\unsure{X|S}}S$ préschéma loc. noethérien \uline{connexe}

$X$ $S$-préschéma \uline{propre} et \uline{simple}

$e$ une section de $X$ sur $S$

$s\in S$ tel que $X_s=f^{-1}(s)$ soit un schéma abélien\deleted{\unsure{ne}} sur $k(s)$.

Alors il existe sur $X$ une seule structure de schéma abélien sur $S$, admettant $e$ comme section unité.

\uline{Cor.} \quad $\overline{\mathrm M}$ si $S$ n'est pas connexe, une structure abélienne est déterminée par la connaissance de sa section unité. (On va \unsure{le|la} \unsure{retrouver|utiliser})

\deleted{\uline{Théorème 2} \quad $S$ \deleted{\unsure{locam|local}} noeth.}

\deleted{$X_1$ et $X_2$ \inserted{de type fini} au-dessus de $S$ ($\varphi_i:X_i\to S$). On suppose}

\deleted{(i) les $\varphi_i$ plats et $\varphi_{2*}(\underline O_{X_2})=\underline O_S$, \uline{ou} $\varphi_2$ séparable}

\deleted{(ii) les fibres de $X_1$ connexes, ou $X_1$ est connexe.}

\deleted{$X_1$ muni d'une section $s_1$}

\deleted{\deleted{Alors \unsure{soit|tout}} $G$ schéma en groupes sur $S$, \deleted{abélien},}
\margin{?}

\deleted{$u:X_1\times_S X_2\longrightarrow G$, \inserted{\unsure{Considérons}} $u_2:X_2\longrightarrow G$,}

\deleted{$\exists\,u_1:X_1\longrightarrow S$ uniques, s.th.}
\[
\deleted{u=\deleted{\unsure{u_1+u_2+0}}u_1\mathrm{pr}_1+u_2\mathrm{pr}_2}
\]

\deleted{\uline{Cor.} \quad If $X_1=X_2$\deleted{\unsure{id. groupe}} \inserted{$=X$} \unsure{group scheme} over $S$, fibres are connected and simple, $X/S$ flat, $u:X\longrightarrow G$ s.th. $u(e)=e$, Then $u$ is a group hom.}

\diagram{\begin{tikzcd}[ampersand replacement=\&]
\deleted{X\times X}\arrow[r]\arrow[d,"\unsure{u{,}u|u\times u}"'] \& \deleted{X}\arrow[d,"u"]\\
G\times G\arrow[r] \& G
\end{tikzcd}}
